\documentclass[tikz,border=8pt]{standalone}
\usepackage{ctex}
\usepackage{amssymb}
\usepackage{pgfplots}
\pgfplotsset{compat=1.18}
\usetikzlibrary{calc,positioning,arrows.meta}

\begin{document}
\begin{tikzpicture}

%% ===== 左图：2D 等高线 + 梯度箭头 =====
\begin{scope}[xshift=0cm]

  % 等高线（同心椭圆，模拟 f(x,y) = x^2 + 2y^2）
  \draw[gray!50, thin] (0,0) ellipse (0.7 and 0.5);
  \draw[gray!50, thin] (0,0) ellipse (1.4 and 1.0);
  \draw[gray!50, thin] (0,0) ellipse (2.1 and 1.5);
  \draw[gray!50, thin] (0,0) ellipse (2.8 and 2.0);

  % 等高线标签
  \node[font=\scriptsize, gray] at (0.5, 0.35) {$c_1$};
  \node[font=\scriptsize, gray] at (1.0, 0.70) {$c_2$};
  \node[font=\scriptsize, gray] at (1.55, 1.05) {$c_3$};

  % 坐标轴
  \draw[->, gray!60, thin] (-3.2,0) -- (3.2,0) node[right, font=\scriptsize] {$x_1$};
  \draw[->, gray!60, thin] (0,-2.3) -- (0,2.3) node[above, font=\scriptsize] {$x_2$};

  % 选取几个点，画梯度箭头（梯度方向 = (2x, 4y)，归一化后缩放）
  % 点 A = (1.2, 0.6)，梯度 (2.4, 2.4)，方向 (1,1)/sqrt(2)
  \coordinate (A) at (1.2, 0.6);
  \draw[->, thick, red!80!black, >=Stealth] (A) -- ++(0.55, 0.55)
    node[above right, font=\scriptsize] {$\nabla f$};
  \fill (A) circle (1.5pt);

  % 点 B = (-1.0, 1.0)，梯度 (-2, 4)，方向 (-1,2)/sqrt(5)
  \coordinate (B) at (-1.0, 1.0);
  \draw[->, thick, red!80!black, >=Stealth] (B) -- ++(-0.27, 0.54)
    node[left, font=\scriptsize] {$\nabla f$};
  \fill (B) circle (1.5pt);

  % 点 C = (1.4, -0.8)，梯度 (2.8, -3.2)，方向 (2.8,-3.2)/norm
  \coordinate (C) at (1.4, -0.8);
  \draw[->, thick, red!80!black, >=Stealth] (C) -- ++(0.42, -0.48)
    node[right, font=\scriptsize] {$\nabla f$};
  \fill (C) circle (1.5pt);

  % 负梯度 / 下降方向 示意（蓝色虚线）
  \draw[->, thick, blue!70!black, dashed, >=Stealth] (A) -- ++(-0.55,-0.55)
    node[below left, font=\scriptsize] {$-\nabla f$};

  % 垂直符号：在 A 点画一个小方块表示梯度垂直等高线
  \draw[gray, thin] ($(A)+(0.1,0)$) -- ($(A)+(0.1,0.1)$) -- ($(A)+(0,0.1)$);

  % 极小值点
  \fill[red!60!black] (0,0) circle (2pt) node[below right, font=\scriptsize] {$\mathbf{x}^*$};

  % 标题
  \node[font=\small\bfseries] at (0, -2.7) {梯度垂直等高线，指向最速上升};

\end{scope}

%% ===== 右图：梯度作为列向量示意 =====
\begin{scope}[xshift=7.5cm]

  % 函数曲面轮廓（简化为抛物线截面）
  \draw[gray!40, thin] (-2.8,0) -- (2.8,0);
  \draw[blue!60!black, thick, domain=-2.5:2.5, samples=60]
    plot (\x, {0.5*\x*\x - 2.0}) node[right, font=\scriptsize] {$f(x_1,x_2)$};

  % 一个具体点 x0 = (1.5, -0.875)
  \coordinate (X0) at (1.5, -0.875);
  \fill (X0) circle (2pt);
  \node[below right, font=\scriptsize] at (X0) {$\mathbf{x}_0$};

  % 切线（斜率 = f'(x_0) = x_0 = 1.5，过 (1.5, -0.875)）
  \draw[red!80!black, thick, domain=0.5:2.5]
    plot (\x, {1.5*(\x-1.5) - 0.875});

  % 梯度箭头（向上指，表示增大方向）
  \draw[->, thick, red!80!black, >=Stealth]
    (X0) -- ++(0.5, 0.75) node[right, font=\scriptsize] {$\nabla f(\mathbf{x}_0)$};

  % 负梯度（向下降方向）
  \draw[->, thick, blue!70!black, dashed, >=Stealth]
    (X0) -- ++(-0.5, -0.75) node[left, font=\scriptsize] {$-\nabla f$};

  % 梯度列向量框
  \node[draw=gray!50, thin, fill=white, rounded corners=2pt,
        font=\small, inner sep=4pt] at (0, 1.8)
    {$\nabla f = (\partial f/\partial x_1,\,\ldots,\,\partial f/\partial x_n)^\top \in \mathbb{R}^n$};

  % 坐标轴
  \draw[->, gray!60, thin] (-3.0, -2.3) -- (3.0, -2.3) node[right, font=\scriptsize] {$x$};
  \draw[->, gray!60, thin] (-3.0, -2.3) -- (-3.0, 2.5) node[above, font=\scriptsize] {$f$};

  % 标题
  \node[font=\small\bfseries] at (0, -2.9) {梯度 $= $ 偏导数列向量};

\end{scope}

%% 整体标题
\node[font=\large\bfseries] at (7.5, 3.3) {梯度几何意义：垂直等高线 + 列向量结构};

\end{tikzpicture}
\end{document}
